% !TeX program = pdflatex
% Self-contained: all example listings are embedded; no external figures/files.
\documentclass[11pt,a4paper]{article}
\usepackage[T1]{fontenc}
\usepackage[utf8]{inputenc}
\usepackage{lmodern}
\usepackage[margin=24mm,headheight=14pt]{geometry}
\usepackage{microtype,amsmath,amssymb,array,booktabs,longtable,enumitem}
\usepackage{xcolor,listings,xurl,hyperref,fancyhdr}
\definecolor{ink}{HTML}{16324F}
\hypersetup{colorlinks=true,linkcolor=ink,urlcolor=ink,pdftitle={UVLM Analysis Guide: Implemented Interfaces and Convergence Studies}}
\newcommand{\code}{\nolinkurl}
\newcolumntype{L}[1]{>{\raggedright\arraybackslash}p{#1}}
\setlength{\parindent}{0pt}
\setlength{\parskip}{.5em}
\setlength{\emergencystretch}{3em}
\renewcommand{\arraystretch}{1.2}
\setlist{leftmargin=*,itemsep=.2em}
\setcounter{tocdepth}{2}
\pagestyle{fancy}\fancyhf{}
\fancyhead[L]{\small UVLM analysis and convergence}\fancyfoot[C]{\thepage}
\lstset{basicstyle=\ttfamily\scriptsize,breaklines=true,columns=fullflexible,
 keepspaces=true,showstringspaces=false,frame=single,rulecolor=\color{black!20},
 backgroundcolor=\color{black!3},tabsize=4,
 literate={Δ}{{$\Delta$}}1 {Γ}{{$\Gamma$}}1 {≈}{{$\approx$}}1 {–}{{-}}1}
\begin{document}
\begin{titlepage}
\vspace*{22mm}
{\Large Implemented terminology and organization migration\par}
\vspace{10mm}
{\Huge\bfseries UVLM analysis guide\par}
\vspace{8mm}
{\Large What each code file does, how data move,\par
and how to define and run an analysis\par}
\vspace{15mm}
\textbf{Delivery date: 15 September 2026.}

This guide describes executable interfaces introduced in the local
\code{WingPropellerUVLM} and \code{WingPropellerUVLMStudies} packages.
It includes the aerodynamic and aeroelastic convergence studies, the retained
Chang linear structural benchmark, and moving-block damping.

The user authorized this terminology and organization migration. The earlier
LaTeX proposal remains the design roadmap; its illustrative bridge constructors
are not an implementation reference. This document identifies the interfaces
that actually exist and explains the remaining architectural work.

\textbf{Structural coupling status.} The retained Chang solver still supplies
the structural equations in the aeroelastic examples. No native AeroBeams
\code{DynamicProblem}--UVLM structural bridge is delivered by this interface
migration. AeroBeams time integration and UVLM wake integration will need the
motion/load mapping and coupled-step coordination described in the roadmap.
\vfill
\textbf{Reading order:} choose an analysis, follow the setup and worked example,
then use the file descriptions and embedded source listings to inspect details.
No earlier proposal or external bibliography is required to understand this guide.
\end{titlepage}
\tableofcontents\clearpage

\section{What changed and what to choose}

Public construction now follows the AeroBeams pattern: create the physical
model, create the analysis problem, call a qualified \code{solve!}, and inspect
results. Public descriptive settings use camelCase. Operation names retain
snake\_case; mutating operations end in \code{!}. Numerical kernels and archived
result keys retain their established names behind a compatibility boundary.

\begin{longtable}{@{}L{.25\linewidth}L{.32\linewidth}L{.35\linewidth}@{}}
\toprule Intended calculation & Implemented entry & Responsibility\\\midrule\endhead
Steady rigid aerodynamics & \code{create_UVLMSteadyProblem} & Existing steady backend, with full surface interaction.\\
Unsteady prescribed geometry & \code{create_UVLMDynamicProblem} & Circulation, dimensional loads and a free wake, with prescribed surface grids.\\
One Chang aeroelastic response & \code{run_chang_linear_aeroelastic.jl} and \code{chang_case.jl} & Retained linear structure, generalized-alpha integration, partitioned UVLM loads.\\
Aerodynamic convergence & \code{chang_convergence_aerodynamic.jl} & Rigid wing/rotor refinement and verified aerodynamic selection.\\
Aeroelastic convergence & \code{chang_convergence_aeroelastic.jl} & Response refinement and moving-block pitch/yaw damping from an aerodynamic selection.\\
Native AeroBeams structure with UVLM & Future bridge, not callable yet & Requires frame-aware kinematics, conservative loads and synchronized step acceptance.\\
\bottomrule
\end{longtable}

The generic aerodynamic facade initially supports the corrected Imperial
dimensional force formulation, a positive fixed core radius and an empty
initial wake. The studies retain their existing fixed/chord/segment core
choices and benchmark-specific rotor geometry. A prescribed RPM stored in an
operating point does not automatically rotate a generic model: its motion
callback must rotate the appropriate grids.

\section{Package organization and setup}
\label{sec:setup}

Paths in this guide are relative to the repository root unless stated
otherwise. Run the shell commands from that root. The aerodynamic library can
be loaded alone; the studies package adds FFT and plotting dependencies.
Plotting code is loaded only when requested, although Plots is declared in the
studies project so enabling it does not depend on the user's global environment.

\begin{lstlisting}
lib/WingPropellerUVLM/
  Project.toml
  src/OperatingPoint.jl
  src/UVLMModel.jl
  src/UVLMProblem.jl
  src/UVLMState.jl                 # existing backend snapshots
  src/backend/                    # retained numerical kernels
  src/wing_propeller/              # grid, rotor and motion helpers
  src/aeroelastic/                 # retained linear time integration
  studies/Project.toml
  studies/Manifest.toml
  studies/setup.jl
  studies/src/WingPropellerUVLMStudies.jl
  studies/src/Convergence.jl
  studies/src/StudySettings.jl
  studies/src/StudyRunner.jl
  studies/src/ChangProblem.jl
  studies/src/Provenance.jl
  studies/src/MovingBlockDamping.jl
  studies/test/runtests.jl
  examples/prescribed_wing.jl
  examples/chang_linear_aeroelastic/run_chang_linear_aeroelastic.jl
  examples/chang_linear_aeroelastic/chang_case.jl
  examples/chang_linear_aeroelastic/studies/convergence/
    chang_convergence_aerodynamic.jl
    chang_convergence_aeroelastic.jl
    src/                          # legacy include forwarding files
\end{lstlisting}

The files under the old convergence \code{src} directory now forward to the
shared implementation. The old moving-block file paths also forward. There
is one maintained implementation of each study component. Existing archived
study scripts remain accessible; their old algorithms and settings do not
become newly validated simply because their includes still work.

For a first installation, use Julia with access to its package registry:
\begin{lstlisting}
julia --project=lib/WingPropellerUVLM/studies lib/WingPropellerUVLM/studies/setup.jl
\end{lstlisting}
The setup script activates the studies project, develops the local aerodynamic
package and instantiates dependencies. The committed study manifest records
the Julia 1.11.5 dependency environment used for this delivery and uses a
relative local package path. Package source files never activate environments.

Start an interactive session with:
\begin{lstlisting}
julia --project=lib/WingPropellerUVLM/studies
\end{lstlisting}
In that REPL, \code{include} of either convergence entry defines a module
and functions without running a simulation. Call \code{main()} explicitly when ready. Direct execution of these
entry files activates the studies project and calls their main function.
Changing a numerical source file after loading a study requires a fresh Julia
process before solving; this prevents a newly hashed source tree from being
attributed to previously loaded numerical code.

\section{Objects and ownership}

\code{UVLMModel} owns copied reference vertex grids, coefficient reference
quantities, symmetry flags and interaction groups. Treat these arrays as
read-only after construction. Changing panel counts means constructing a new
model. Two problems can use the same model because each allocates its own
backend state and workspace.

\code{OperatingPoint} holds airspeed, density, incidence, sideslip and
prescribed rotor-speed metadata. \code{UVLMSolver} holds the fixed core radius,
interaction switch and wake-shedding fraction. The operating point is the
physical condition; the solver describes the selected aerodynamic algorithm.

\code{UVLMState} contains the accepted backend system, \code{timeNow} and
\code{stepIndex}. \code{UVLMWorkspace} contains a separate trial system and
flags indicating whether a trial is open and ready for acceptance. Each
system still uses the existing backend \code{System} layout. This is a
transitional facade, not a complete field-by-field backend storage rewrite.
Deep copies isolate trials and rollback. That choice prioritizes correctness
and is more expensive than a future reusable snapshot-buffer implementation.

\code{UVLMResults} owns saved arrays. Its circulation and geometry records
are copied, so subsequent solves cannot overwrite past samples through
workspace aliases. \code{UVLMDynamicProblem} brings these objects together
with the time grid, motion callback and history controls.

The Chang facade has a different result contract. Its existing \code{ChangModel}
stores resolved benchmark settings and matrices; \code{ChangDynamicProblem}
calls the existing run orchestration and exposes its returned workspace,
solution and postprocessed results. It adds \code{completed} and
\code{solverConverged}. Its backend fields and CSV keys retain legacy spelling.

\section{Units, names and compatibility}

\begin{longtable}{@{}L{.30\linewidth}L{.29\linewidth}L{.33\linewidth}@{}}
\toprule Existing term & New public term & Meaning\\\midrule\endhead
\code{speed_mps} & \code{airspeed} & m/s; used by operating points and convergence settings.\\
\code{end_time_s} & \code{finalTime} & Seconds.\\
\code{maximum_iterations} & \code{maximumIterations} & Outer coupling iteration limit.\\
\code{state_tolerance}, \code{load_tolerance} & \code{stateTolerance}, \code{loadTolerance} & Same existing residual definitions.\\
\code{trim_revolutions} & \code{baselineRevolutions} & Startup load preparation, not a deformed structural trim.\\
\code{wing_span}, \code{wing_chord} & \code{wingSpanPanels}, \code{wingChordPanels} & Study refinement controls.\\
\code{prop_radial}, \code{prop_chord} & \code{propellerRadialPanels}, \code{propellerChordPanels} & Panels per blade and direction.\\
\code{spanwise_panels}, \code{chordwise_panels} & \code{nSpanwisePanels}, \code{nChordwisePanels} & Benchmark geometry settings.\\
\code{azimuth_deg} & \code{azimuthStepDeg} & Study azimuth increment in degrees.\\
\code{core_radius_m} & \code{coreRadius} & Fixed radius in metres.\\
\code{make_plots}, \code{dry_run} & \code{makePlots}, \code{dryRun} & Output and execution controls.\\
\code{minimum_fit_r_squared} & \code{minimumFitRSquared} & Dimensionless fit-quality threshold.\\
\bottomrule
\end{longtable}

The generic operating-point fields \code{angleOfAttack} and \code{sideslip}
are radians. The migrated editable benchmark/study settings retain explicit
unit suffixes: \code{alphaDeg}, \code{angleOfAttackDeg} and \code{azimuthStepDeg}.
The single-case legacy key \code{azimuth_step_deg} and study key
\code{azimuth_deg} share the public spelling \code{azimuthStepDeg}.
Their separate boundaries preserve the appropriate legacy key: the Chang
constructor uses \code{legacy_case_settings}, and studies use \code{legacy_settings}.
Do not strip a \code{Deg} suffix and retain the same number. Other benchmark
names use systematic camelCase conversion, such as \code{freestreamSpeedMps},
\code{rotationRpm} and \code{pitchStiffnessNmPerRad}. They are still explicit
physical inputs. This initial interface uses safe unit-preserving aliases;
it does not implement every speculative name from the proposal.

\code{canonical_settings} converts a legacy study NamedTuple into public
settings. \code{legacy_settings} converts a study back for existing runners.
Duplicate aliases in one input are rejected. Study constructors accept a
base settings NamedTuple and recursively merge keyword overrides, rejecting
unknown override keys. Enum values such as \code{:fixed}, \code{:record_fraction}
and \code{:exact_virtual_work} retain their spelling. Refinement-family symbols
are translated with the corresponding control names.

\section{Worked example: prescribed wing in the time domain}
\label{sec:wing}

The complete executable \code{examples/prescribed_wing.jl} is reproduced in
the appendix. It defines a one-metre-span, one-metre-chord wing with four
spanwise and two chordwise panels. The reference area, chord and span are
1 in SI units. Airspeed is 10 m/s and incidence is 3 degrees, explicitly
converted to radians. The example prescribes a small vertical displacement:
\begin{equation}
 z(t)=0.002\sin(2\pi\,2t)\quad\text{metres}.
\end{equation}
The callback copies reference grids and adds this displacement. Geometry
differences between successive time levels supply the backend surface-motion
terms. The wing is prescribed; no mass or stiffness is inferred from its grids.

\begin{lstlisting}
import WingPropellerUVLM
include("lib/WingPropellerUVLM/examples/prescribed_wing.jl")
problem = PrescribedWingExample.create_problem()
WingPropellerUVLM.solve!(problem)
times = problem.results.savedTimeVector
forces = problem.results.forceOverTime
moments = problem.results.momentOverTime
\end{lstlisting}
The example's \code{main()} is also a complete create-and-solve call.
This short run demonstrates the interface; it is not a converged aerodynamic
benchmark or a damping record.

Grid indexing is \code{(coordinate,chordwise vertex,spanwise vertex)}. Thus a
surface with $n_c$ by $n_s$ panels has shape $(3,n_c+1,n_s+1)$. A model accepts
a vector of grids, one for each aerodynamic surface. All motion callbacks must
return the same surface count and topology with finite coordinates.

For a wing followed by four blades of one rotor, a suitable interaction-group
vector is \code{[1,2,2,2,2]}. With \code{interaction=true}, the global solve
includes cross-group influence. With \code{interaction=false}, groups are
isolated aerodynamically, while the four blades in group 2 still interact.
For a second rotor, assign a third group to all its blades. Symmetry flags are
specified per surface; a symmetric wing and explicit physical blades use
\code{[true,false,false,false,false]}. In the Chang aeroelastic analysis,
disabling aerodynamic cross-group influence does not remove the structural
wing--propeller connection.

For a stationary wing, omit \code{surfaceMotion}; the default returns the
reference grids. For a rotating blade, rotate that blade's reference grid about
its shaft and include the hub offset. Use $\psi(t)=\psi_0+\Omega t$ with the
chosen spin sign. Setting the operating point's speed alone does not define
the shaft axis, hub position or blade geometry.

\subsection{Time and wake settings}
Supply either \code{timeVector} or \code{initialTime}, the Julia keyword
$\Delta t$, and \code{finalTime}. The latter creates the range
\code{initialTime:} $\Delta t$ \code{:finalTime}; if the interval is not an
integer number of steps, the final stored grid coordinate lies below the
requested bound. Completion means completion of that resolved discrete grid.
An explicit vector must contain at least two finite, strictly increasing times.
It may be prescribed nonuniformly, but the moving-block FFT requires uniform
response samples. No adaptive error controller is introduced here.

\code{maximumWakeRows} is an integer per surface, or one integer applied to
all surfaces. Capacity is positive; active rows start at zero and increase
once per accepted step, up to the allocated limit. Retaining ten rows at a
constant 0.005-s step retains approximately 0.05 s of recent wake history.
Changing the time step changes that interpretation.

\subsection{History and force interpretation}
\code{trackingTimeSteps=true} enables saved histories.
\code{trackingFrequency=1} saves every accepted step; larger values save every
$k$th step and the final step. The initial empty-wake state is not a solved
force sample and is not appended to these histories. \code{saveGeometry=true}
additionally stores surface-panel geometry, with the same stride.

\code{forceOverTime} contains total dimensional vertex forces in N;
\code{momentOverTime} contains moments in N m about \code{model.reference.r}.
They use backend coordinates. These totals sum the explicitly modeled
surfaces. A symmetry image contributes aerodynamic influence but is not a
second physical surface summed into these arrays. Do not infer a full-wing
result by doubling arbitrary force or moment components without considering
reflection signs. Use \code{dimensional_loads(problem)} for individual vertex
forces and their matching application positions.

\section{A physical step and its numerical trials}

For accepted aerodynamic history $A_n$, an attempted next time step evaluates
geometry at $t_{n+1}$ and solves
\begin{equation}
 A(\mathcal G_{n+1},\mathcal W_n)\Gamma_{n+1}=b(\mathcal G_{n+1},\mathcal W_n,v).
\end{equation}
Repeated trials share the accepted history $\mathcal W_n$. They may have
different proposed surface geometry. Only acceptance convects and sheds the
wake for the next physical step.

\begin{lstlisting}
UVLM.begin_time_step!(problem)
loads = UVLM.evaluate_trial!(problem)
# Optional further trials at this same physical time:
# loads = UVLM.evaluate_trial!(problem; surfaces=trialGrids)
UVLM.commit_time_step!(problem)
# Or discard the attempted step with UVLM.rollback_time_step!(problem).
\end{lstlisting}
Here \code{UVLM} is an imported alias for \code{WingPropellerUVLM}. Opening a
second step while one is open is an error. Evaluating without opening a step,
committing without a successful trial, and committing twice are errors.
Rollback discards trial changes and does not append a result sample.

\code{solve!} repeatedly opens, evaluates and commits until the grid is
complete. On an evaluation or wake-convection error it rolls back the attempted
step, stores the error as \code{terminationReason}, and rethrows it. Accepted
earlier samples remain available. Fix the cause before resuming or construct
a new problem. A failed calculation must not be presented as a complete result.

This API supplies aerodynamic transactions for future coupling. It does not
provide a structural trial or a structural residual. The future bridge must
also snapshot AeroBeams state, map rotations and velocities, transfer loads
conservatively, converge the structural/aerodynamic pair and accept both together.

\section{Worked example: one Chang aeroelastic analysis}

\code{examples/chang_linear_aeroelastic/chang_case.jl} contains the editable
single-case defaults. The entry \code{run_chang_linear_aeroelastic.jl} loads
those settings and applies \code{CHANG_*} environment overrides. When included,
call \code{run_chang_example()} explicitly.
The convergence studies define their own settings; editing the single-case
airspeed does not change a study.

The studies API can also construct a problem from the same defaults, with
explicit overrides and no environment overrides:
\begin{lstlisting}
import WingPropellerUVLMStudies as Studies
include("lib/WingPropellerUVLM/examples/chang_linear_aeroelastic/chang_case.jl")
settings = Studies.merge_settings(Studies.canonical_settings(chang_case_defaults()),
    (; simulation=(; freestreamSpeedMps=85.0, finalTime=6.0),
       output=(; plotResults=false, animateWake=false)))
model = Studies.create_ChangModel(settings)
problem = Studies.create_ChangDynamicProblem(;model)
Studies.solve!(problem)
problem.results.completed
problem.results.solverConverged
solution = problem.results.solution
\end{lstlisting}

The example settings are grouped by responsibility:
\begin{longtable}{@{}L{.20\linewidth}L{.72\linewidth}@{}}
\toprule Group & Inputs and effect\\\midrule\endhead
\code{wing} & Root/tip chord, span, symmetry and panel counts. Spanwise panels still determine the retained structural element count.\\
\code{propeller} & Radius, chord, blade count, radial/chordwise panels, installation fractions, reference RPM/speed and collective offset.\\
\code{simulation} & Density, freestream speed, angles, azimuth increment, final time, interaction switch, force model and moment projection.\\
\code{aerodynamic} & Fixed core or factor law, elastic-axis fraction and hub-load arm.\\
\code{wake} & Wing wake row capacity and rotor retained revolutions, with optional explicit row overrides.\\
\code{excitation} & Baseline averaging and the pitch impulse's time, duration and magnitude.\\
\code{integration} & Generalized-alpha numerical dissipation and state/amplitude limits.\\
\code{coupling} & Outer iterations, residual tolerances and relaxation.\\
\code{output} & Directory/label, response plots, animation and frame stride.\\
\code{structural} & Physical damping settings, pylon/rotor modal properties, mass and inertias derived from the configured data.\\
\bottomrule
\end{longtable}

The structural equations remain $M\ddot q+C\dot q+Kq=Q$. Generalized-alpha
advances this linear benchmark; the aerodynamic load still depends on the
trial geometry and wake. The outer partitioned iterations make the structural
state and aerodynamic load consistent before wake acceptance.
\code{rhoInf} controls numerical integration dissipation. It is not a physical
modal damping ratio. Physical damping inputs remain in the structural group.

\code{solution.displacement_history}, \code{velocity_history} and
\code{acceleration_history} hold the retained structural coordinates.
\code{last_step} reports accepted steps; coupling iteration and residual
arrays describe convergence. The facade preserves the benchmark's existing
output writer and file schemas. It does not rename Euler coordinates to
AeroBeams Wiener--Milenkovic parameters or reinterpret the benchmark vector
as an AeroBeams global state.

Startup baseline averaging subtracts the mean load about the configured
geometry. It is not a solved deformed equilibrium. For damping interpretation,
choose a response interval after baseline preparation and excitation, and
inspect whether periodic rotor forcing or multiple modes dominate that interval.

\section{Aerodynamic convergence: editing and executing}
\label{sec:aerostudy}

Both existing convergence entry filenames are preserved. Their
\code{create_study()} functions now return typed study objects. The legacy
\code{aerodynamic_study()} and \code{aeroelastic_study()} accessors still
return snake\_case NamedTuples for existing scripts.

\begin{lstlisting}
import WingPropellerUVLMStudies as Studies
include("lib/WingPropellerUVLM/examples/chang_linear_aeroelastic/studies/convergence/chang_convergence_aerodynamic.jl")
base = ChangConvergenceAerodynamic.create_study()
study = Studies.create_AerodynamicConvergenceStudy(base.settings;
    dryRun=true, makePlots=false,
    outputDirectory="output/my_aerodynamic_matrix")
Studies.solve!(study)
study.results.directory
\end{lstlisting}

The first run above writes the case matrix only. Review it before setting
\code{dryRun=false} for the same settings in a fresh problem. Every enabled
family needs at least three strictly ordered levels. Panel counts and retained
wake revolutions increase with refinement; core radius and azimuth increment
decrease. An azimuth step must divide 360 degrees and be no larger than
90 degrees. These are validity limits, not recommended accuracy settings.

\code{physical} supplies the rigid geometry and operating condition. Its
\code{rpm} is the actual RPM at its \code{airspeed}. Changing one through a
partial override does not silently change the other. To preserve the advance
ratio, explicitly set $N_{new}=N_{ref}U_{new}/U_{ref}$ when defining this stage.
The editable defaults show this relationship directly.

\code{nominal} holds the controls used when another family is varied.
\code{sweeps} holds the level lists and \code{families} chooses which lists
are active. \code{coreMode=:fixed} interprets \code{core} in metres;
\code{:chord} interprets it as a fraction of full local chord;
\code{:segment} interprets it as a fraction of the local span/radial segment.
These are different regularization laws and must not be compared as if their
numbers had the same units.

\code{simulatedRevolutions} is the total rigid run length.
\code{averagedRevolutions} is the end interval for periodic metrics.
\code{absoluteTolerances}, \code{relativeTolerance} and
\code{periodicTolerances} test CT/CQ changes and periodic histories.
\code{reuseExisting} permits reuse only when metadata, source and files match.
\code{confirmSelection} requests the combined candidate/finer verification
after individual refinement families pass.

\textbf{Interpretation limit.} Existing aerodynamic acceptance is based on
propeller CT/CQ, including periodic comparisons. Wing lift is recorded but
does not currently gate that selection. With disabled aerodynamic interaction,
CT/CQ cannot establish wing-load convergence. This migration preserves that
criterion; do not describe the resulting selection as proof of every wing
force and moment being converged. Adding wing metrics is a separate study-design
change identified in the roadmap.

\section{Aeroelastic convergence and the selection handoff}

First obtain \code{aerodynamic_selection.toml} from a successful aerodynamic
stage under the current source identity. A matrix-only run does not create a
verified selection. The loader checks the selection and its recorded evidence;
an unverified hand-written file is not an equivalent input.

\begin{lstlisting}
import WingPropellerUVLMStudies as Studies
include("lib/WingPropellerUVLM/examples/chang_linear_aeroelastic/studies/convergence/chang_convergence_aeroelastic.jl")
base = ChangConvergenceAeroelastic.create_study()
study = Studies.create_AeroelasticConvergenceStudy(base.settings;
    selectionFile="output/my_aerodynamic_run/aerodynamic_selection.toml",
    outputDirectory="output/my_aeroelastic_run",
    airspeed=85.0, dryRun=true, makePlots=false,
    response=(;finalTime=6.0))
Studies.solve!(study)
\end{lstlisting}

Here the top-level \code{airspeed=nothing} would retain the selected operating
point. A positive override changes only the aeroelastic operating point and
scales its prescribed rotor RPM by $U_{new}/U_{reference}$. The original
aerodynamic reference is kept in the result and report. Its seal is not
relabelled as a verification at the overridden speed.

At fixed azimuth increment $\Delta\psi$ and angular speed $\Omega$,
$\Delta t=\Delta\psi/|\Omega|$ with radians in this equation. For a positive
RPM $N$ and a degree-valued study increment, $\Delta t=\Delta\psi_{deg}/(6N)$.
Changing speed therefore changes physical time resolution even when the
azimuth resolution is unchanged.

The aeroelastic \code{sweeps} can contain \code{nothing}. The runner then
constructs three levels starting at the selected aerodynamic control and
refining it. Explicit lists must also begin at the selected value. The wing
spanwise family is excluded because this benchmark still ties it to the
structural element count. Independent aerodynamic and structural meshes are
future bridge work, not an implemented capability of this retained model.

\code{response} defines the simulation end, baseline averaging and pitch
excitation. \code{structural}, \code{integration} and \code{coupling} retain
their meanings from the single-case analysis. \code{damping} selects moving-block
settings and study comparison tolerances. An aeroelastic dry run still needs a
valid selection and validates the translated structural configurations.

\section{Moving-block damping and what the reported numbers mean}

The shared \code{MovingBlockDamping.jl} preserves the selected reference
algorithm: positive peaks bound a fitting segment; block length occupies the
configured fraction of that segment; blocks shift one sample; each block has
its mean removed; a rectangular FFT identifies the dominant non-DC component;
and a line is fitted to log amplitude versus block-start time.

For $x(t)\simeq A_0e^{\lambda t}\sin(2\pi f t+\varphi)$,
\begin{equation}
 \log A_k\simeq a+\lambda t_k,\qquad
 \omega_d=2\pi f,\qquad
 \zeta=-\lambda/\sqrt{\lambda^2+\omega_d^2}.
\end{equation}
Negative growth rate $\lambda$ (s$^{-1}$) means decay; positive means growth.
$f$ is Hz; $\omega_d$ is rad/s; $\zeta$ is dimensionless. A percentage is
$100\zeta$. Existing serialized fields such as
\code{moving_block_lambda_per_s} and \code{damping_ratio_percent} remain
explicit; the migration does not reinterpret percentages as fractions.

\begin{longtable}{@{}L{.33\linewidth}L{.59\linewidth}@{}}
\toprule Setting & Effect\\\midrule\endhead
\code{fitStartS} & Explicit start, or automatic post-excitation start when nothing.\\
\code{blockSizeMode} & \code{:record_fraction} preserves reference sizing; \code{:duration} requests a fixed-duration window.\\
\code{blockSize}, \code{sizeRatioLb}, \code{sizeRatioUb} & Reference sample count and permitted lower/upper segment fractions; defaults 512, 0.25, 0.5.\\
\code{peakFromStart}, \code{peakFromEnd} & Select the positive peaks bounding the fitting segment.\\
\code{blockDurationS}, \code{blockOverlap} & Used in duration mode; do not control the reference one-sample shift mode.\\
\code{frequencyBandHz} & Allowed frequency interval for acceptance.\\
\code{minimumPeaks}, \code{minimumFitRSquared} & Minimum signal/fit evidence.\\
\code{lambdaTolerancePerS}, \code{frequencyToleranceHz} & Differences allowed between refinement results.\\
\bottomrule
\end{longtable}

Uniform response sampling is required. Saving animation frames sparsely does
not thin the lightweight response CSV used by this estimator. FFT bin spacing
is $1/(N_b\Delta t)$; a tolerance finer than this spacing does not itself
establish frequency accuracy. Inspect the raw response, block frequency history,
log-amplitude fit and fit-quality diagnostics. Positive-peak screening remains
sensitive to offsets, and a periodically forced or multimode response need not
yield a free-mode damping ratio merely because a fit is available.

\section{Results, failures, archives and provenance}

\begin{longtable}{@{}L{.37\linewidth}L{.55\linewidth}@{}}
\toprule File or object & How to use it\\\midrule\endhead
\code{case_matrix.csv} & Exact requested controls, time increment and wake rows; generated even for dry runs.\\
\code{study.toml}, \code{report.md} & Study settings, method identity and interpretation.\\
Rigid case history and metadata & CT/CQ and wing-load histories plus the source/options/file checks used before reuse.\\
\code{aerodynamic_selection.toml} & Verified handoff following individual and combined aerodynamic checks.\\
\code{aerodynamic_input.toml} & Aeroelastic stage's selected aerodynamic input and reference record.\\
\code{resolved_configuration.toml}, \code{run.log} & Actual translated Chang settings and solver trace for one case.\\
\code{analysis_status.toml} & Exact accepted/expected step counts and coupling completion for an aeroelastic case.\\
\code{response_history.csv} & Lightweight accepted response used for damping.\\
\code{damping_fit.toml} & Pitch/yaw fit summaries, estimator identity and validity.\\
\code{pitch_moving_block.csv}, \code{yaw_moving_block.csv} & Block-start times, log amplitudes, fitted values and tracked frequencies.\\
\code{study.results} & Existing structured case rows, directory and selected/reference data after solving.\\
\bottomrule
\end{longtable}

Finishing a process, converging every physical step, obtaining a valid damping
fit, and passing a refinement criterion are distinct events. The migrated
aeroelastic runner checks accepted step count and coupling flags before fitting.
Cases stopped by limits cannot pass solely because their final sample happens
to lie close to the requested end time.

\code{Provenance.jl} hashes the aerodynamic source, shared study source, Chang
source, project/manifest files and Julia version. The identity deliberately
covers indirect dependencies. The estimator also records its own hash.
This delivery uses one conservative numerical source identity; independently
minimal aerodynamic/structural/coupling dependency fingerprints remain future
work. The version marker for this migration is \code{uvlm-migration-v2}.

File moves change source identity even when numerical formulas do not change.
Existing selection files and seals are retained untouched. Consequently an
old selection can fail current-source verification. Use
\code{Studies.inspect_selection(path)} to read its archived TOML record;
the returned \code{verifiedForCurrentSource=false} is intentional. Rerun the
required aerodynamic evidence under the migrated source to obtain a new
verified selection. Do not edit a hash to force acceptance.

\section{Code-by-code responsibility map}
\label{sec:files}

The following tables cover the active numerical library, benchmark and study
files. Each row states where to investigate a behavior or make a future
change. Paths in the first table are relative to \code{lib/WingPropellerUVLM/src}.

\begin{longtable}{@{}L{.36\linewidth}L{.56\linewidth}@{}}
\toprule File & Responsibility\\\midrule\endhead
\code{WingPropellerUVLM.jl} & Includes numerical components and exports both new constructors and compatibility operations.\\
\code{OperatingPoint.jl} & Validates flow inputs; resolves prescribed or constant-advance-ratio rotor metadata; constructs the backend Freestream.\\
\code{UVLMModel.jl} & Validates/copied grids and references; defines UVLMSolver; translates the camelCase partitioned constructor to existing solver options.\\
\code{UVLMProblem.jl} & Defines accepted/trial/results containers, dynamic/steady constructors, trial geometry updates, stepping, rollback and dimensional result recording.\\
\code{UVLMState.jl} & Existing backend snapshots, restore operations, wake-row bookkeeping and noncommitting trial support used by the retained benchmark.\\
\code{backend/system.jl} & Backend array allocation and storage layout. Model, history and scratch fields still coexist here behind the new facade.\\
\code{backend/panel.jl} & Surface and wake panel representations and panel-level geometry operations.\\
\code{backend/geometry.jl} & Converts grid vertices into panels, collocation/force geometry and associated geometric data.\\
\code{backend/reference.jl}, \code{freestream.jl} & Reference area/chord/span/density and coordinate-frame/freestream conventions.\\
\code{backend/induced.jl} & Vortex-segment induced velocities and influence assembly, including finite-core and interaction handling.\\
\code{backend/circulation.jl} & Boundary-condition right-hand side and circulation-system solution.\\
\code{backend/analyses.jl} & Steady and unsteady orchestration, one-step propagation and separate wake advancement.\\
\code{backend/wake.jl} & Wake velocities, convection and shedding. Wake history is chronological numerical state.\\
\code{backend/nearfield.jl} & Corrected dimensional Imperial segment loads; selected by the new generic facade.\\
\code{backend/nearfield_postprocessing.jl} & Nodal load assembly, matching application points and coefficient/result reduction.\\
\code{backend/legacy_nearfield.jl} & Retained alternative force implementation for explicit legacy comparisons.\\
\code{backend/farfield.jl}, \code{stability.jl} & Existing far-field and stability/postprocessing utilities; no new native coupled eigenproblem is introduced.\\
\code{backend/nonlinear.jl} & Retained nonlinear aerodynamic routines outside the initial generic facade's supported path.\\
\code{backend/rotors.jl}, \code{vspgeom.jl} & Existing rotor/import utilities. Generic model construction initially receives already defined vertex grids.\\
\code{backend/visualization.jl} & Existing aerodynamic geometry/solution visualization export.\\
\code{wing_propeller/BladeGeometry.jl} & Tabulated/interpolated blade geometry and twist definitions.\\
\code{wing_propeller/GridUtilities.jl} & Interpolation and wing/propeller grid construction.\\
\code{wing_propeller/Initialization.jl} & Existing case-oriented combined wing/rotor backend initialization.\\
\code{wing_propeller/Kinematics.jl} & Updates grid positions and rotor/deformation motion.\\
\code{aeroelastic/GeneralizedAlpha.jl} & Generalized-alpha parameters/corrector, scaled partitioned residuals and iteration.\\
\code{aeroelastic/Excitations.jl} & Smooth pulse shape and generalized excitation loads.\\
\bottomrule
\end{longtable}

The next table uses paths relative to \code{studies/src}.
\begin{longtable}{@{}L{.35\linewidth}L{.57\linewidth}@{}}
\toprule File & Responsibility\\\midrule\endhead
\code{WingPropellerUVLMStudies.jl} & Package entry, shared module inclusion, public exports and loaded source identity.\\
\code{StudySettings.jl} & Canonical/legacy key translation, duplicate rejection, recursive overrides, typed study constructors, solve dispatch and archival inspection.\\
\code{StudyRunner.jl} & Verifies the aerodynamic selection, applies the aeroelastic speed/RPM override, invokes the central moving-block run and writes its operating-point report.\\
\code{ChangProblem.jl} & Builds the retained Chang model from explicit settings and wraps its run as create/solve/result operations.\\
\code{Convergence.jl} & Case definitions, ordered sweeps, case reuse/execution, refinement comparisons, plots, tables and study orchestration.\\
\code{aerodynamic_options.jl} & Existing rigid aerodynamic option definitions, validity checks and per-case metadata/file integrity.\\
\code{aerodynamic_backend.jl} & Rigid combined wing/rotor calculation and dimensional load reduction.\\
\code{aerodynamic_selection.jl} & Selects controls, requires combined verification, writes and verifies the handoff evidence.\\
\code{aeroelastic_adapter.jl} & Converts selected aerodynamic controls to explicit Chang settings, launches cases, checks grid completion and collects fits.\\
\code{MovingBlockDamping.jl} & Reference moving-block estimator, screening, block FFT, growth/frequency fit and diagnostics.\\
\code{moving_block_adapter.jl} & Connects the chosen estimator to convergence validation/output; stores its identity and diagnostics.\\
\code{damping_metrics.jl} & Existing response reader and legacy estimator retained for compatibility and archived workflows. Current typed aeroelastic studies use MovingBlockDamping.\\
\code{Provenance.jl} & Conservative numerical source identity including indirect study and benchmark dependencies.\\
\bottomrule
\end{longtable}

The retained benchmark sources live in
\code{examples/chang_linear_aeroelastic/src}.
\begin{longtable}{@{}L{.37\linewidth}L{.55\linewidth}@{}}
\toprule File & Responsibility\\\midrule\endhead
\code{ChangAeroelastic.jl} & Assembles the case module, defines ChangModel, builds workspaces and coordinates solving/postprocessing. The new wrapper passes its already built model into this orchestration.\\
\code{chang_configuration.jl} & Parses/resolves legacy benchmark settings, automatic defaults and units; validates physical/numerical choices. New constructors explicitly disable environment overrides.\\
\code{chang_model_parameters.jl} & Resolves model dimensions, attachments, rotor speed, time grid, reference quantities and derived parameters.\\
\code{chang_structural_matrices.jl} & Beam, pylon, mass/inertia and gyroscopic matrix primitives.\\
\code{chang_structural_model.jl} & Assembles the retained benchmark's constrained M/C/K system and structural metadata.\\
\code{chang_workspaces.jl} & Allocates combined UVLM/structural working arrays and wake/motion context.\\
\code{chang_uvlm_coupling.jl} & Trial structural-to-grid motion and aerodynamic generalized loads, including the benchmark's virtual-work moment projection.\\
\code{chang_simulation.jl} & Accepted time-step loop, baseline preparation, excitation, partitioned iteration, wake commit and failure/limit handling.\\
\code{chang_postprocessing.jl} & Extracts physical response histories from structural coordinates and writes requested output.\\
\code{chang_plotting.jl}, \code{chang_animation.jl} & Optional time-history plots and sampled geometry/wake animation.\\
\code{chang_propeller_trim.jl} & Retained specialized propeller/baseline helper; its name does not turn startup load subtraction into a native AeroBeams trim solution.\\
\bottomrule
\end{longtable}

\section{Public function contracts and call paths}

\begin{longtable}{@{}L{.36\linewidth}L{.56\linewidth}@{}}
\toprule Function & Contract\\\midrule\endhead
\code{create_OperatingPoint} & Requires positive finite airspeed/density. Prescribed RPM or reference-RPM scaling is resolved once; conflicting policies are rejected.\\
\code{create_UVLMModel} & Requires finite vertex grids and positive reference quantities; copies input geometry and validates per-surface flags/groups.\\
\code{create_UVLMSolver} & Requires a positive fixed core and shedding fraction between zero and one.\\
\code{create_PartitionedCoupling} & Returns the existing validated PartitionedCouplingOptions. Its backend fields retain snake\_case.\\
\code{create_UVLMDynamicProblem} & Resolves time grid, initializes geometry and empty wake, and allocates independent accepted/trial storage.\\
\code{begin_time_step!} & Opens the next scheduled interval without advancing accepted history.\\
\code{evaluate_trial!} & Restores a copy of accepted state, applies trial geometry, solves circulation/loads and returns dimensional forces/positions.\\
\code{commit_time_step!} & Convects/sheds on a private copy, installs accepted history, updates counts/time and saves requested output.\\
\code{rollback_time_step!} & Discards trial storage and restores attempted-step flags to accepted time.\\
\code{create_UVLMSteadyProblem}, \code{solve!} & Wrap the existing full-interaction steady backend; solved backend results are in \code{problem.system}.\\
\code{create_ChangModel}, \code{create_ChangDynamicProblem} & Resolve explicit benchmark inputs, construct retained structural data and describe one run.\\
\code{create_AerodynamicConvergenceStudy}, \code{create_AeroelasticConvergenceStudy} & Copy/normalize settings, apply validated overrides and create an unsolved study object.\\
\code{study_settings}, \code{legacy_settings} & Return copied canonical settings or their compatibility representation.\\
\code{Studies.solve!} & Runs the selected study/benchmark, records results and preserves existing evidence rules.\\
\bottomrule
\end{longtable}

The aerodynamic call path is: entry settings $\to$ typed study $\to$
Convergence case matrix $\to$ rigid backend $\to$ periodic CT/CQ histories
$\to$ individual comparisons $\to$ combined verification $\to$ selection.

The aeroelastic path is: verified selection $\to$ optional speed/RPM override
$\to$ translated Chang configuration $\to$ structural model/workspace
$\to$ partitioned accepted steps $\to$ complete response history
$\to$ moving-block pitch/yaw fits $\to$ refinement comparisons.

The generic prescribed-motion path is: grid model + operating point + solver
$\to$ dynamic problem $\to$ begin/evaluate/commit $\to$ dimensional histories.
This third path does not call the Chang structural matrices or the convergence
runner.

\section{Verification and remaining migration work}

Before edits, the existing suite passed 434 assertions. The original numerical
fingerprint is recorded in \code{docs/migration/baseline.toml}. Migration tests
compare old one-call propagation to the new trial/commit path for circulation,
temporal rates, forces and wake geometry, including a mature wake and capacity
truncation. They check A--B--A trials, failed geometry, rollback, duplicate
commit, copied results, steady parity, settings translation and dry-run behavior.
The existing moving-block reference regressions and selection tests are retained.

\textbf{Final result: 508 assertions passed in the declared studies environment
on Julia 1.11.5.} This comprises the 434 existing checks and 74 additional checks.
A 38-step Chang response at 85 m/s has exactly matching displacement, velocity
and coupling-iteration histories through the legacy and new interfaces. A
deliberately stopped case is rejected before damping acceptance. The lockfile
was resolved with Julia Pkg and checked against the declared compatibility
bounds, including Interpolations 0.15.1. These checks validate the migration;
they do not establish a full six-second damping response or a converged
production aerodynamic selection.

Run the full library and study suite with:
\begin{lstlisting}
julia --project=lib/WingPropellerUVLM/studies lib/WingPropellerUVLM/studies/test/runtests.jl
\end{lstlisting}
That runner restricts the package load path to the declared project and Julia
standard libraries. The core library suite is separated from study/FFT tests
so the aerodynamic package can also be tested in isolation. Validation results
and any environment limits are recorded in \code{docs/migration/VALIDATION.md}.

This delivery introduces working facade objects and migrates shared studies.
It does not claim completion of every item in the P0--P11 roadmap. Remaining
work includes a deeper backend field-ownership split and buffer optimization,
fully separated writers for the retained Chang path, minimal stage-specific
dependency fingerprints, additional wing-load convergence criteria and
independent structural/aerodynamic meshes. The native AeroBeams bridge,
external structural load/step hooks, rotor structural correspondence, adaptive
coupling and eigen/periodic stability remain separate numerical work.

\section{Compiling this self-contained document}
All listings below are embedded text. They do not execute Julia during LaTeX
compilation, and no external source files, images or bibliography are required.
Run twice to resolve the contents and references:
\begin{lstlisting}
pdflatex -interaction=nonstopmode -halt-on-error UVLM_ANALYSIS_GUIDE.tex
pdflatex -interaction=nonstopmode -halt-on-error UVLM_ANALYSIS_GUIDE.tex
\end{lstlisting}
The source can also be uploaded alone to Overleaf with pdfLaTeX selected.
The authoring environment did not provide a LaTeX compiler; rendered PDF
inspection is therefore not asserted by the source-level checks.

\appendix
\section{Complete executable entries and facade sources}
These listings reproduce the delivered code. Comments describe the retained physical conventions; the main text explains how each entry is used. File paths identify the source location, but compilation does not read those files.
\subsection{\texorpdfstring{\code{prescribed_wing.jl}}{prescribed wing.jl}}
\code{lib/WingPropellerUVLM/examples/prescribed_wing.jl}
\begin{lstlisting}
# Run with --project=lib/WingPropellerUVLM. Including defines the module only.
module PrescribedWingExample
import WingPropellerUVLM as UVLM

function create_problem()
    grid,_=UVLM.wing_to_grid([0.,0.],[0.,1.],[0.,0.],ones(2),zeros(2),zeros(2),4,2)
    model=UVLM.create_UVLMModel(;surfaces=[grid],
        reference=UVLM.Reference(1.,1.,1.,zeros(3),10.),name="Prescribed heaving wing")
    point=UVLM.create_OperatingPoint(;airspeed=10.,density=1.225,angleOfAttack=deg2rad(3.))
    solver=UVLM.create_UVLMSolver(;coreRadius=.001)
    function motion(model,t)
        grids=deepcopy(model.grids)
        grids[1][3,:,:] .+= .002*sin(2pi*2t)
        return grids
    end
    return UVLM.create_UVLMDynamicProblem(;model,operatingPoint=point,aeroSolver=solver,
        timeVector=collect(0.:.005:.05),maximumWakeRows=10,surfaceMotion=motion,
        trackingTimeSteps=true,trackingFrequency=1,saveGeometry=false)
end
main()=UVLM.solve!(create_problem())
end
if abspath(PROGRAM_FILE)==@__FILE__
    problem=PrescribedWingExample.main()
    println("Completed: ",problem.results.completed)
    println("Last force [N]: ",last(problem.results.forceOverTime))
end

\end{lstlisting}
\subsection{\texorpdfstring{\code{chang_convergence_aerodynamic.jl}}{chang convergence aerodynamic.jl}}
\code{lib/WingPropellerUVLM/examples/chang_linear_aeroelastic/studies/convergence/chang_convergence_aerodynamic.jl}
\begin{lstlisting}
# Aerodynamic convergence. Edit camelCase settings in create_study().
# Including this file defines the entry only. Use the studies project in the REPL.
if abspath(PROGRAM_FILE) == @__FILE__
    import Pkg
    Pkg.activate(normpath(joinpath(@__DIR__,"..","..","..","..","studies")))
end

module ChangConvergenceAerodynamic
import WingPropellerUVLMStudies
const Convergence=WingPropellerUVLMStudies.Convergence

function create_study()
    airspeed = 65.0
    trimSpeed = 65.0
    trimRpm = 1212.0
    propellerRadius = 1.15
    muProp = trimSpeed / (trimRpm * 2 * pi / 60.0) / propellerRadius
    omega = airspeed / propellerRadius / muProp
    rpm = omega * 60.0 / (2 * pi)
    physical = (
        airspeed = airspeed,
        densityKgpm3 = 1.225,
        alphaDeg = 3.0,
        betaDeg = 0.0,
        wingSpanM = 7.5,
        wingRootChordM = 1.8,
        wingTipChordM = 1.8,
        wingSymmetric = true, # Wing/wake images only. Blades never use symmetry.
        elasticAxisFraction = 0.30,
        propellerRadiusM = propellerRadius,
        propellerChordM = 0.197,
        blades = 4,
        attachmentEta = 0.83,
        pylonLengthM = 5.6 * 0.3048,
        rpm = rpm, # Actual RPM at airspeed; derived from the constant mu_prop.
        collectiveOffsetDeg = 0.0,
        interaction = false,
        wakeSheddingFraction = 0.1,
    )

    # Core mode: :fixed (metres), :chord (fraction of full local chord),
    # or :segment (fraction of the local 3D span/radial vortex edge).
    coreMode = :fixed
    nominal = (wingSpanPanels=30, wingChordPanels=10, propellerRadialPanels=10, propellerChordPanels=10,
        wakeRevolutions=2.0, core=0.001, azimuthStepDeg=5.0)

    # Absolute panel counts per direction (propeller counts are per blade).
    # Each family varies one parameter; others retain the nominal values.
    # Use at least three ordered levels for every enabled family.
    sweeps = (
        wingSpanPanels = [20, 40, 60],
        wingChordPanels = [5, 10, 20],
        propellerRadialPanels = [5, 10, 20],
        propellerChordPanels = [5, 10, 20],
        wakeRevolutions = [1.0, 2.0, 3.0],
        core = [0.01, 0.003, 0.001, 0.0003],
        azimuthStepDeg = [5.0, 2.5, 1.25],
    )
    settings = (;
        physical, coreMode, nominal, sweeps,
        families = collect(keys(sweeps)),
        simulatedRevolutions = 8,
        averagedRevolutions = 2,
        absoluteTolerances = (CT=1e-6, CQ=1e-6),
        relativeTolerance = 0.01,
        periodicTolerances = (CT=1e-6, CQ=1e-6),
        dryRun = false, # true writes only the case matrix.
        makePlots = true,
        reuseExisting = true,
        # Publish a selection only after every family AND a combined
        # candidate/finer verification satisfy the CT/CQ criteria.
        confirmSelection = true,
        outputDirectory = normpath(joinpath(@__DIR__,"..","..","output","convergence_aerodynamic")),
    )
    return WingPropellerUVLMStudies.create_AerodynamicConvergenceStudy(settings)
end

# Legacy settings access retains the original schema for existing scripts.
aerodynamic_study()=WingPropellerUVLMStudies.legacy_settings(create_study())
main()=WingPropellerUVLMStudies.solve!(create_study()).results
end

if abspath(PROGRAM_FILE) == @__FILE__
    ChangConvergenceAerodynamic.main()
end

\end{lstlisting}
\subsection{\texorpdfstring{\code{chang_convergence_aeroelastic.jl}}{chang convergence aeroelastic.jl}}
\code{lib/WingPropellerUVLM/examples/chang_linear_aeroelastic/studies/convergence/chang_convergence_aeroelastic.jl}
\begin{lstlisting}
# Aeroelastic convergence. Edit camelCase settings in create_study().
# Including this file defines the entry only. Use the studies project in the REPL.
if abspath(PROGRAM_FILE) == @__FILE__
    import Pkg
    Pkg.activate(normpath(joinpath(@__DIR__,"..","..","..","..","studies")))
end

module ChangConvergenceAeroelastic
import WingPropellerUVLMStudies
const Convergence=WingPropellerUVLMStudies.Convergence
const MovingBlockStudy=WingPropellerUVLMStudies.MovingBlockStudy
const MovingBlockDamping=WingPropellerUVLMStudies.MovingBlockDamping
operating_selection(selection,s)=WingPropellerUVLMStudies.operating_selection(selection,s)
run_aeroelastic(s=aeroelastic_study())=WingPropellerUVLMStudies.run_aeroelastic(s)

function create_study()
    settings = (;
        selectionFile = normpath(joinpath(@__DIR__,"..","..","output",
            "convergence_aerodynamic","aerodynamic_selection.toml")),
        outputDirectory = normpath(joinpath(@__DIR__,"..","..","output","convergence_aeroelastic")),
        dryRun = false,
        makePlots = true,

        # nothing: aerodynamic speed. Set e.g. 85.0 for this aeroelastic study only.
        # RPM scales by new_speed / aerodynamic_speed (constant advance ratio).
        airspeed = 80,

        # nothing: selected aerodynamic value followed by two finer values.
        # Supply explicit absolute lists to choose your own refinements.
        # Lists must start at the selected value and have >=3 ordered levels.
        # Wing span remains fixed at the selected structural element count.
        sweeps = (
            wingChordPanels = nothing,    # e.g. [10, 12, 15] if 10 was selected
            propellerRadialPanels = nothing,
            propellerChordPanels = nothing,
            wakeRevolutions = nothing,
            core = nothing,
            azimuthStepDeg = nothing,
        ),
        families = [:wingChordPanels,:propellerRadialPanels,:propellerChordPanels,:wakeRevolutions,:core,:azimuthStepDeg],

        response = (
            finalTime = 6.0,
            baselineRevolutions = 10.0,
            baselineAverageRevolutions = 1.0,
            impulseStartS = nothing, # automatic: after trim_revolutions
            impulseDurationS = 0.15,
            impulseMagnitudeNm = 1200.0,
        ),
        damping = (
            fitStartS = nothing, # after pulse + one revolution; at least 0.75 s
            # moving_block_2.jl: rectangular, mean-removed FFT blocks shifted
            # one sample; lambda is the slope of log amplitude against time.
            blockSizeMode = :record_fraction, # or :duration for a fixed window
            blockSize = 512, # doubled/halved to satisfy the ratios below
            sizeRatioLb = 0.25,
            sizeRatioUb = 0.50,
            peakFromStart = 1,
            peakFromEnd = 0,
            blockDurationS = 1.0, # used only with :duration
            blockOverlap = 0.9,    # used only with :duration
            frequencyBandHz = (3.0,6.0), # acceptance check on GLOBAL dominant frequency
            minimumPeaks = 8,
            minimumFitRSquared = 0.8,
            lambdaTolerancePerS = 0.01,
            frequencyToleranceHz = 0.1,
        ),
        structural = (
            propellerDampingRatio = 0.0,
            stiffnessDampingRatio = 0.0,
            dampingReferenceFrequencyHz = nothing,
            pitchFrequencyHz = 7.97,
            yawFrequencyHz = 7.97,
            pitchStiffnessNmPerRad = 19220.0,
            yawStiffnessNmPerRad = 18916.0,
            twistFrequencyHz = 12.73,
            twistStiffnessNmPerRad = 16835.0,
            pylonMassPerLengthKgpm = 0.0506 * 14.5939029372064 / 0.3048,
            bladeMassKg = 1.44,
        ),
        integration = (rhoInf=1.0, stateNormLimit=1e3, propellerAngleLimitDeg=Inf),
        coupling = (maximumIterations=10,stateTolerance=1e-5,loadTolerance=1e-2,
            equilibriumTolerance=1e-10,coupledEquilibriumTolerance=1e-4,relaxation=1.0),
        propellerMomentProjection = :exact_virtual_work,
        hubLoadArmFactor = 0.5,
    )
    return WingPropellerUVLMStudies.create_AeroelasticConvergenceStudy(settings)
end

# Legacy settings access retains the original schema for existing scripts.
aeroelastic_study()=WingPropellerUVLMStudies.legacy_settings(create_study())
main()=WingPropellerUVLMStudies.solve!(create_study()).results
end

if abspath(PROGRAM_FILE) == @__FILE__
    ChangConvergenceAeroelastic.main()
end

\end{lstlisting}
\subsection{\texorpdfstring{\code{OperatingPoint.jl}}{OperatingPoint.jl}}
\code{lib/WingPropellerUVLM/src/OperatingPoint.jl}
\begin{lstlisting}
"""Resolved SI flow and prescribed rotor operating point. Angles are radians."""
struct OperatingPoint
    airspeed::Float64
    density::Float64
    angleOfAttack::Float64
    sideslip::Float64
    angularSpeed::Float64
    referenceAirspeed::Float64
    referenceRotationRPM::Float64
    rotorSpeedPolicy::Symbol
end

function create_OperatingPoint(; airspeed, density=1.225, angleOfAttack=0.0,
    sideslip=0.0, rotationRPM=nothing, referenceAirspeed=airspeed,
    referenceRotationRPM=nothing, rotorSpeedPolicy=:prescribed)
    all(x -> x isa Real && !(x isa Bool) && isfinite(x),
        (airspeed,density,angleOfAttack,sideslip,referenceAirspeed)) ||
        throw(ArgumentError("Operating point must contain finite real values"))
    airspeed>0 && density>0 && referenceAirspeed>0 ||
        throw(ArgumentError("Airspeed, reference airspeed and density must be positive"))
    rotorSpeedPolicy in (:prescribed,:constantAdvanceRatio) ||
        throw(ArgumentError("rotorSpeedPolicy must be :prescribed or :constantAdvanceRatio"))
    if rotorSpeedPolicy==:constantAdvanceRatio
        isnothing(rotationRPM) || throw(ArgumentError("Specify referenceRotationRPM, not rotationRPM, for proportional scaling"))
        isnothing(referenceRotationRPM) && throw(ArgumentError("referenceRotationRPM is required"))
        rpm=referenceRotationRPM*airspeed/referenceAirspeed
    else
        isnothing(referenceRotationRPM) || throw(ArgumentError("Use rotationRPM for :prescribed speed"))
        rpm=something(rotationRPM,0.0)
        referenceRotationRPM=rpm
    end
    all(x -> x isa Real && !(x isa Bool) && isfinite(x),(rpm,referenceRotationRPM)) ||
        throw(ArgumentError("Rotor speeds must be finite"))
    return OperatingPoint(airspeed,density,angleOfAttack,sideslip,2pi*rpm/60,
        referenceAirspeed,referenceRotationRPM,rotorSpeedPolicy)
end
freestream(p::OperatingPoint) = Freestream(p.airspeed,p.angleOfAttack,p.sideslip,zeros(3))


\end{lstlisting}
\subsection{\texorpdfstring{\code{UVLMModel.jl}}{UVLMModel.jl}}
\code{lib/WingPropellerUVLM/src/UVLMModel.jl}
\begin{lstlisting}
"""Reference aerodynamic vertex grids and coefficient references; treat as read-only.
Grids have size (3, nChordwisePanels+1, nSpanwisePanels+1), in metres.
"""
struct UVLMModel{G,R}
    name::String
    grids::G
    reference::R
    symmetric::Vector{Bool}
    interactionGroups::Vector{Int}
end
function create_UVLMModel(; surfaces,reference,name="",symmetric=false,interactionGroups=nothing)
    isempty(surfaces) && throw(ArgumentError("At least one vertex grid is required"))
    for g in surfaces
        ndims(g)==3 && size(g,1)==3 && size(g,2)>=2 && size(g,3)>=2 ||
            throw(ArgumentError("surfaces must contain (3,nChord+1,nSpan+1) vertex grids"))
        all(isfinite,g) || throw(ArgumentError("Geometry must be finite"))
    end
    reference isa Reference || throw(ArgumentError("reference must be a Reference"))
    all(x -> isfinite(x) && x>0,(reference.S,reference.c,reference.b,reference.V,reference.rho)) &&
        all(isfinite,reference.r) || throw(ArgumentError("Invalid reference quantities"))
    n=length(surfaces)
    sym=symmetric isa Bool ? fill(symmetric,n) : Bool.(symmetric)
    ids=isnothing(interactionGroups) ? collect(1:n) : Int.(interactionGroups)
    length(sym)==length(ids)==n || throw(DimensionMismatch("One symmetry/group entry per surface is required"))
    return UVLMModel(String(name),[Float64.(g) for g in surfaces],deepcopy(reference),sym,ids)
end

"""Aerodynamic algorithm controls. AIC assembly remains enabled for moving geometry."""
struct UVLMSolver
    coreRadius::Float64
    interaction::Bool
    wakeSheddingFraction::Float64
end
function create_UVLMSolver(;coreRadius=1e-3,interaction=true,wakeSheddingFraction=.1)
    isfinite(coreRadius) && coreRadius>0 || throw(ArgumentError("coreRadius must be positive"))
    isfinite(wakeSheddingFraction) && 0<=wakeSheddingFraction<=1 ||
        throw(ArgumentError("wakeSheddingFraction must lie in [0,1]"))
    return UVLMSolver(coreRadius,interaction,wakeSheddingFraction)
end

"""AeroBeams-style constructor for the existing linear partitioned solver options."""
function create_PartitionedCoupling(;maximumIterations=10,stateTolerance=1e-5,
    loadTolerance=1e-2,equilibriumTolerance=1e-10,coupledEquilibriumTolerance=1e-4,
    relaxation=1.0)
    options=PartitionedCouplingOptions(;maximum_iterations=maximumIterations,
        state_tolerance=stateTolerance,load_tolerance=loadTolerance,
        equilibrium_tolerance=equilibriumTolerance,
        coupled_equilibrium_tolerance=coupledEquilibriumTolerance,relaxation)
    validate_partitioned_coupling_options(options)
    return options
end


\end{lstlisting}
\subsection{\texorpdfstring{\code{UVLMProblem.jl}}{UVLMProblem.jl}}
\code{lib/WingPropellerUVLM/src/UVLMProblem.jl}
\begin{lstlisting}
"""Accepted aerodynamic state. The backend System remains a transitional storage facade."""
mutable struct UVLMState{S}
    system::S
    timeNow::Float64
    stepIndex::Int
end
"""Private trial System and scratch storage, isolated from the accepted state."""
mutable struct UVLMWorkspace{S}
    system::S
    trialReady::Bool
    stepOpen::Bool
    timeEndTimeStep::Float64
end
mutable struct UVLMResults
    completed::Bool
    terminationReason::String
    savedTimeVector::Vector{Float64}
    circulationOverTime::Vector{Vector{Float64}}
    forceOverTime::Vector{Vector{Float64}}
    momentOverTime::Vector{Vector{Float64}}
    surfaceGeometryOverTime::Vector{Any}
end
UVLMResults()=UVLMResults(false,"initialized",Float64[],Vector{Float64}[],
    Vector{Float64}[],Vector{Float64}[],Any[])

mutable struct UVLMDynamicProblem{M,F,S}
    model::M
    aeroSolver::UVLMSolver
    operatingPoint::OperatingPoint
    timeVector::Vector{Float64}
    surfaceMotion::F
    state::UVLMState{S}
    workspace::UVLMWorkspace{S}
    trackingTimeSteps::Bool
    trackingFrequency::Int
    saveGeometry::Bool
    results::UVLMResults
end

function initialize_backend(model,solver,point,maximumWakeRows,grids)
    length(grids)==length(model.grids) || throw(DimensionMismatch("Surface count changed"))
    system=System(deepcopy(model.grids);nw=maximumWakeRows)
    system.reference[]=Reference(model.reference.S,model.reference.c,model.reference.b,
        model.reference.r,point.airspeed,point.density)
    system.freestream[]=freestream(point)
    system.symmetric .= model.symmetric
    system.surface_id .= model.interactionGroups
    apply_trial_geometry!(system,grids,solver)
    copy_surfaces_to_previous!(system,length(system.surfaces))
    return system
end
function apply_trial_geometry!(system,grids,solver)
    length(grids)==length(system.grids) || throw(DimensionMismatch("Surface count changed"))
    for i in eachindex(grids)
        size(grids[i])==size(system.grids[i]) || throw(DimensionMismatch("Grid topology changed"))
        all(isfinite,grids[i]) || throw(ArgumentError("Nonfinite trial geometry"))
        system.grids[i] .= grids[i]
        _,ratio,panels=grid_to_surface_panels(system.grids[i];fcore=(c,ds)->solver.coreRadius)
        system.ratios[i] .= ratio
        system.surfaces[i] .= panels
    end
    return system
end

"""Create a fixed-grid or prescribed-motion aerodynamic analysis.
`surfaceMotion(model,time)` returns all vertex grids at that time, without
mutating model. No structural equations or automatic rotor motion are inferred.
Specify either timeVector or initialTime/Δt/finalTime. Initial wake is empty.
"""
function create_UVLMDynamicProblem(;model,aeroSolver=create_UVLMSolver(),operatingPoint,
    initialTime=0.0,Δt=nothing,finalTime=nothing,timeVector=nothing,
    maximumWakeRows=100,surfaceMotion=(m,t)->m.grids,
    trackingTimeSteps=true,trackingFrequency=1,saveGeometry=false)
    trackingFrequency isa Integer && trackingFrequency>0 || throw(ArgumentError("trackingFrequency must be positive"))
    if isnothing(timeVector)
        !isnothing(Δt) && !isnothing(finalTime) || throw(ArgumentError("Provide Δt and finalTime, or timeVector"))
        all(isfinite,(initialTime,Δt,finalTime)) && Δt>0 && finalTime>initialTime ||
            throw(ArgumentError("Invalid time interval"))
        times=collect(Float64(initialTime):Float64(Δt):Float64(finalTime))
    else
        isnothing(Δt) && isnothing(finalTime) || throw(ArgumentError("timeVector conflicts with Δt/finalTime"))
        times=Float64.(timeVector)
    end
    length(times)>=2 && all(isfinite,times) && all(>(0),diff(times)) ||
        throw(ArgumentError("At least two finite, increasing time samples are required"))
    n=length(model.grids)
    rows=maximumWakeRows isa Integer ? fill(Int(maximumWakeRows),n) : Int.(maximumWakeRows)
    length(rows)==n && all(>(0),rows) || throw(ArgumentError("Positive wake capacity required for every surface"))
    sys=initialize_backend(model,aeroSolver,operatingPoint,rows,surfaceMotion(model,first(times)))
    state=UVLMState(sys,first(times),1)
    workspace=UVLMWorkspace(deepcopy(sys),false,false,first(times))
    return UVLMDynamicProblem(model,aeroSolver,operatingPoint,times,surfaceMotion,
        state,workspace,trackingTimeSteps,trackingFrequency,saveGeometry,UVLMResults())
end

function begin_time_step!(p::UVLMDynamicProblem)
    p.workspace.stepOpen && throw(ArgumentError("A time step is already open"))
    p.state.stepIndex<length(p.timeVector) || throw(ArgumentError("No remaining time steps"))
    p.workspace.timeEndTimeStep=p.timeVector[p.state.stepIndex+1]
    p.workspace.stepOpen=true
    p.workspace.trialReady=false
    return p
end

"""Evaluate a trial from the accepted state. This never commits wake history.
Returns dimensional vertex forces and their matching application points.
"""
function evaluate_trial!(p::UVLMDynamicProblem;surfaces=nothing)
    w=p.workspace
    w.stepOpen || throw(ArgumentError("Call begin_time_step! first"))
    w.trialReady=false
    w.system=deepcopy(p.state.system)
    sys=w.system
    copy_surfaces_to_previous!(sys,length(sys.surfaces))
    grids=isnothing(surfaces) ? p.surfaceMotion(p.model,w.timeEndTimeStep) : surfaces
    apply_trial_geometry!(sys,grids,p.aeroSolver)
    propagate_system!(sys,freestream(p.operatingPoint),w.timeEndTimeStep-p.state.timeNow;
        additional_velocity=nothing,repeated_points=repeated_trailing_edge_points(sys.surfaces),
        nwake=copy(p.state.system.nwake),eta=p.aeroSolver.wakeSheddingFraction,
        calculate_influence_matrix=true,near_field_analysis=true,derivatives=false,
        interaction_id=p.model.interactionGroups,interaction=p.aeroSolver.interaction,advance_wake=false)
    all(isfinite,sys.Γ) && all(isfinite,sys.dΓdt) || error("Nonfinite aerodynamic trial")
    loads=dimensional_loads(sys)
    all(f -> all(v->all(isfinite,v),f),loads.forces) || error("Nonfinite aerodynamic forces")
    w.trialReady=true
    return loads
end

"""Dimensional vertex loads in backend coordinates; positions use the same force geometry."""
function dimensional_loads(system::System)
    return (;forces=imperial_nodal_forces(system),positions=imperial_nodal_positions(system))
end
dimensional_loads(p::UVLMDynamicProblem)=dimensional_loads(p.state.system)

function save_time_step!(p::UVLMDynamicProblem)
    loads=dimensional_loads(p)
    force=zeros(3); moment=zeros(3)
    for i in eachindex(loads.forces), j in eachindex(loads.forces[i])
        f=loads.forces[i][j]; r=loads.positions[i][j]-p.model.reference.r
        force .+= f; moment .+= cross(r,f)
    end
    result=p.results
    push!(result.savedTimeVector,p.state.timeNow)
    push!(result.circulationOverTime,copy(p.state.system.Γ))
    push!(result.forceOverTime,force); push!(result.momentOverTime,moment)
    p.saveGeometry && push!(result.surfaceGeometryOverTime,deepcopy(p.state.system.surfaces))
    return p
end

function commit_time_step!(p::UVLMDynamicProblem)
    w=p.workspace
    w.stepOpen && w.trialReady || throw(ArgumentError("A successful uncommitted trial is required"))
    # Commit on a private copy: a wake-convection failure leaves the trial retryable.
    sys=deepcopy(w.system)
    rows=copy(p.state.system.nwake)
    advance_wake!(sys,freestream(p.operatingPoint),w.timeEndTimeStep-p.state.timeNow;
        nwake=rows,interaction_id=p.model.interactionGroups,interaction=p.aeroSolver.interaction)
    commit_wake_rows!(rows,size.(sys.wakes,1)); sys.nwake .= rows
    p.state.system=sys
    p.state.timeNow=w.timeEndTimeStep; p.state.stepIndex+=1
    w.stepOpen=false; w.trialReady=false
    if p.trackingTimeSteps && (mod(p.state.stepIndex-1,p.trackingFrequency)==0 || p.state.stepIndex==length(p.timeVector))
        save_time_step!(p)
    end
    p.results.completed=p.state.stepIndex==length(p.timeVector)
    p.results.terminationReason=p.results.completed ? "completed" : "running"
    return p
end
function rollback_time_step!(p::UVLMDynamicProblem)
    p.workspace.system=deepcopy(p.state.system)
    p.workspace.stepOpen=false; p.workspace.trialReady=false
    p.workspace.timeEndTimeStep=p.state.timeNow
    return p
end
function solve!(p::UVLMDynamicProblem)
    p.workspace.stepOpen && throw(ArgumentError("Finish or roll back the open trial before solve!"))
    try
        while p.state.stepIndex<length(p.timeVector)
            begin_time_step!(p); evaluate_trial!(p); commit_time_step!(p)
        end
    catch e
        rollback_time_step!(p)
        p.results.terminationReason=sprint(showerror,e)
        rethrow()
    end
    return p
end

mutable struct UVLMSteadyProblem{M}
    model::M
    aeroSolver::UVLMSolver
    operatingPoint::OperatingPoint
    system::Union{Nothing,System}
end
function create_UVLMSteadyProblem(;model,aeroSolver=create_UVLMSolver(),operatingPoint)
    aeroSolver.interaction || throw(ArgumentError("Steady facade currently requires full interaction"))
    return UVLMSteadyProblem(model,aeroSolver,operatingPoint,nothing)
end
function solve!(p::UVLMSteadyProblem)
    m=p.model; op=p.operatingPoint
    ref=Reference(m.reference.S,m.reference.c,m.reference.b,m.reference.r,op.airspeed,op.density)
    p.system=steady_analysis(deepcopy(m.grids),ref,freestream(op);symmetric=m.symmetric,
        surface_id=m.interactionGroups,fcore=(c,ds)->p.aeroSolver.coreRadius,derivatives=false)
    return p
end

\end{lstlisting}
\subsection{\texorpdfstring{\code{StudySettings.jl}}{StudySettings.jl}}
\code{lib/WingPropellerUVLM/studies/src/StudySettings.jl}
\begin{lstlisting}
# Public settings use camelCase; the backend and existing file schemas retain
# their established names/units. This is the single translation boundary.
const SETTING_NAMES=Dict(
    :speed_mps=>:airspeed, :end_time_s=>:finalTime,
    :trim_revolutions=>:baselineRevolutions,
    :trim_average_revolutions=>:baselineAverageRevolutions,
    :spanwise_panels=>:nSpanwisePanels, :chordwise_panels=>:nChordwisePanels,
    :radial_panels=>:nRadialPanels, :core_radius_m=>:coreRadius,
    :wing_span=>:wingSpanPanels, :wing_chord=>:wingChordPanels,
    :prop_radial=>:propellerRadialPanels, :prop_chord=>:propellerChordPanels,
    :azimuth_deg=>:azimuthStepDeg)
const LEGACY_NAMES=Dict(v=>k for (k,v) in SETTING_NAMES)
function canonical_key(k::Symbol)
    haskey(SETTING_NAMES,k) && return SETTING_NAMES[k]
    parts=split(string(k),'_')
    return Symbol(first(parts)*join(uppercasefirst.(parts[2:end])))
end
function legacy_key(k::Symbol)
    haskey(LEGACY_NAMES,k) && return LEGACY_NAMES[k]
    k in (:CT,:CQ) && return k
    separated=replace(string(k),r"([A-Z])([A-Z][a-z])"=>s"\1_\2")
    return Symbol(lowercase(replace(separated,r"([a-z0-9])([A-Z])"=>s"\1_\2")))
end
function translate_settings(x::NamedTuple,keyfn)
    pairsOut=Pair{Symbol,Any}[]
    seen=Set{Symbol}()
    for (key,value) in pairs(x)
        mapped=keyfn(key)
        mapped in seen && throw(ArgumentError("Duplicate aliases for $mapped"))
        push!(seen,mapped)
        converted=key==:families ? keyfn.(value) : translate_settings(value,keyfn)
        push!(pairsOut,mapped=>converted)
    end
    return (;pairsOut...)
end
translate_settings(x,keyfn)=deepcopy(x)
canonical_settings(x::NamedTuple)=translate_settings(x,canonical_key)
legacy_settings(x::NamedTuple)=translate_settings(x,legacy_key)
function merge_settings(base::NamedTuple,overrides::NamedTuple)
    unknown=setdiff(keys(overrides),keys(base))
    isempty(unknown) || throw(ArgumentError("Unknown settings: $unknown"))
    return (; (k=>(hasproperty(overrides,k) ?
        (v isa NamedTuple && overrides[k] isa NamedTuple ? merge_settings(v,overrides[k]) : deepcopy(overrides[k])) : deepcopy(v))
        for (k,v) in pairs(base))...)
end

abstract type ConvergenceStudy end
mutable struct AerodynamicConvergenceStudy{S} <: ConvergenceStudy
    settings::S
    results::Union{Nothing,NamedTuple}
end
mutable struct AeroelasticConvergenceStudy{S} <: ConvergenceStudy
    settings::S
    results::Union{Nothing,NamedTuple}
end
study_settings(s::ConvergenceStudy)=deepcopy(s.settings)
legacy_settings(s::ConvergenceStudy)=legacy_settings(s.settings)

function create_AerodynamicConvergenceStudy(settings::NamedTuple;kwargs...)
    resolved=merge_settings(canonical_settings(settings),canonical_settings((;kwargs...)))
    Convergence.validate_aerodynamic(legacy_settings(resolved))
    return AerodynamicConvergenceStudy(resolved,nothing)
end
create_AerodynamicConvergenceStudy(;kwargs...)=create_AerodynamicConvergenceStudy((;kwargs...))
function create_AeroelasticConvergenceStudy(settings::NamedTuple;kwargs...)
    resolved=merge_settings(canonical_settings(settings),canonical_settings((;kwargs...)))
    legacy=legacy_settings(resolved)
    hasproperty(legacy,:selection_file) || throw(ArgumentError("selectionFile is required"))
    MovingBlockDamping.moving_block_metrics(Float64[],Float64[];moving_block_options(legacy.damping)...)
    return AeroelasticConvergenceStudy(resolved,nothing)
end
create_AeroelasticConvergenceStudy(;kwargs...)=create_AeroelasticConvergenceStudy((;kwargs...))

function check_loaded_source()
    Convergence.SOURCE()==LOADED_SOURCE || error("Numerical source changed after this module was loaded; start a fresh Julia process")
end
function solve!(s::AerodynamicConvergenceStudy)
    check_loaded_source()
    settings=legacy_settings(s)
    settings.make_plots && (@eval using Plots)
    s.results=Base.invokelatest(Convergence.run_aerodynamic,settings)
    return s
end
function solve!(s::AeroelasticConvergenceStudy)
    check_loaded_source()
    settings=legacy_settings(s)
    settings.make_plots && (@eval using Plots)
    s.results=run_aeroelastic(settings)
    return s
end

"""Read an archived selection without asserting current numerical verification.
The returned record cannot be used as a verified selection by solve!.
"""
function inspect_selection(path)
    payload=Convergence.TOML.parsefile(path)
    return (;verifiedForCurrentSource=false,record=payload)
end

\end{lstlisting}
\subsection{\texorpdfstring{\code{ChangProblem.jl}}{ChangProblem.jl}}
\code{lib/WingPropellerUVLM/studies/src/ChangProblem.jl}
\begin{lstlisting}
"""Build the existing Chang benchmark from explicit camelCase or legacy settings.
No CHANG_* environment overrides are applied by this constructor.
"""
function create_ChangModel(settings::NamedTuple)
    Convergence.load_model()
    moduleRef=Convergence.Model.ChangAeroelastic
    config=Base.invokelatest(moduleRef.load_chang_configuration,legacy_case_settings(settings);env=Dict{String,String}())
    return Base.invokelatest(moduleRef.build_chang_model,config)
end
legacy_case_settings(settings::NamedTuple)=translate_settings(settings,
    key -> key==:azimuthStepDeg ? :azimuth_step_deg : legacy_key(key))
mutable struct ChangDynamicProblem{M}
    model::M
    results::Union{Nothing,NamedTuple}
end
create_ChangDynamicProblem(;model)=ChangDynamicProblem(model,nothing)
function solve!(p::ChangDynamicProblem)
    check_loaded_source()
    p.results===nothing || throw(ArgumentError("Create a fresh ChangDynamicProblem for another run"))
    m=Convergence.Model.ChangAeroelastic
    Base.invokelatest(m.load_chang_visualization,p.model.config.output)
    run=Base.invokelatest(m.run_chang_case,p.model.config;model=p.model)
    completed=run.solution.last_step==length(p.model.parameters.dt)
    solverConverged=completed && all(run.solution.coupling_converged)
    p.results=merge(run,(;completed,solverConverged))
    return p
end

\end{lstlisting}
\subsection{\texorpdfstring{\code{StudyRunner.jl}}{StudyRunner.jl}}
\code{lib/WingPropellerUVLM/studies/src/StudyRunner.jl}
\begin{lstlisting}
function operating_selection(selection,s)
    speed=get(s,:speed_mps,nothing)
    isnothing(speed) && return selection
    speed isa Real && !(speed isa Bool) && isfinite(speed) && speed>0 ||
        throw(ArgumentError("Aeroelastic speed_mps must be finite and positive, or nothing"))
    p=selection.physical
    rpm=p.rpm*(speed/p.speed_mps)
    isfinite(rpm) && rpm>0 || throw(ArgumentError("Scaled aeroelastic RPM must be finite and positive"))
    # The verification seal belongs to the original operating point only.
    return (;physical=merge(p,(;speed_mps=Float64(speed),rpm)),
        core_mode=selection.core_mode,selected=selection.selected,
        aerodynamic_reference=selection)
end

"""Run damping convergence with an optional speed override after verifying the source selection."""
function run_aeroelastic(s)
    check_loaded_source()
    reference=Convergence.load_selection(s.selection_file)
    selection=operating_selection(reference,s)
    println("Aeroelastic operating point: $(selection.physical.speed_mps) m/s, $(selection.physical.rpm) RPM")
    Convergence.load_model()
    result=Base.invokelatest(Convergence.run_aeroelastic_loaded,MovingBlockStudy(s),selection)
    description="Aerodynamic reference: $(reference.physical.speed_mps) m/s, $(reference.physical.rpm) RPM.\n\nAeroelastic operating point: $(selection.physical.speed_mps) m/s, $(selection.physical.rpm) RPM.\n\n"
    if selection.physical.speed_mps!=reference.physical.speed_mps
        description*="The aerodynamic selection was verified at the reference speed; this study checks aeroelastic sensitivity at the overridden speed with proportional RPM.\n\n"
    end
    open(joinpath(result.directory,"report.md"),"a") do io
        print(io,"\n",description)
    end
    return merge(result,(;aerodynamic_reference=reference))
end

\end{lstlisting}

\end{document}
